
\section{Related Work}

\subsection{Low-Rank Compression}

\textbf{SVD factorization} (Eckart \& Young, 1936). The optimal rank-$r$ 
approximation in Frobenius norm.

\textbf{ALBERT} (Lan et al., 2020). Cross-layer parameter sharing.

\textbf{Tensor train} (Novikov et al., 2015). Higher-order factorization.

\textbf{ORION uses SVD} with energy-based rank selection.

\subsection{Linear Attention}

\textbf{Linear Transformer} (Katharopoulos et al., 2020). $\phi(x) = \text{elu}(x) + 1$.

\textbf{Performer} (Choromanski et al., 2021). Random Fourier Features.

\textbf{Cosformer} (Qin et al., 2022). Cosine reweighting.

\textbf{ORION uses RFF} for the kernel heads.

\subsection{Hybrid Attention}

\textbf{Longformer} (Beltagy et al., 2020). Sliding-window + global.

\textbf{BigBird} (Zaheer et al., 2020). Sliding-window + random + global.

\textbf{NASC} (this framework). Three-branch cascade.

\textbf{ORION bifurcates by head}, not by branch -- different heads use different 
mechanisms.

\subsection{Softpick}

\textbf{Alman \& Song (2024).} Sparse alternative to softmax.

\textbf{NASC's softpick} uses ReLU with fallback.

\textbf{ORION's softpick} uses $\text{ReLU}(e^x - 1)$, simpler.

\subsection{Position of Our Work}

\begin{table}[H]
\centering
\small
\begin{tabular}{lll}
\toprule
\textbf{Axis} & \textbf{Prior work} & \textbf{ORION} \\
\midrule
\textbf{Factorization} & Post-hoc SVD & Zero-shot SVD \\
\textbf{Head mechanism} & Uniform & Bifurcated (2/3, 1/3) \\
\textbf{Parameter reduction} & Often none & \textbf{~6$$\times$} \\
\textbf{Robustness} & Varies & \textbf{\texttt{copy\_weights} raises} \\
\bottomrule
\end{tabular}
\end{table}

\textbf{Key distinction.} ORION is one of the few modules in our framework 
that \textbf{actually reduces parameters}. It combines three orthogonal 
techniques, though one of them (the \texttt{copy\_weights} raise) is a bug.



\subsection{Comparison Table}

\begin{lstlisting}

---
\end{lstlisting}